\question[4]
Bestimme, welchen Wert die Variable \texttt{index} nach Ausführung des folgenden Codes hat,
 und begründe deine Antwort.
\begin{lstlisting}
a = [5, 3, 8, 2, 9, 1]
index = -1
for i in range(len(a)):
    if a[i] > 6 and index == -1:
        index = i
print(index)
\end{lstlisting}

\begin{solutionbox}{4cm}

Die Schleife durchsucht die Liste von vorne nach hinten. Sobald
ein Element groesser als 6 gefunden wird und index noch -1 ist
(also noch kein Treffer gespeichert wurde), wird der Index
gemerkt. Da a[2]=8 das erste Element groesser als 6 ist, wird
index auf 2 gesetzt und bleibt wegen der Bedingung index == -1
danach unverändert.

Ausgabe: 2

\end{solutionbox}